\chapter{Fundamentação Teórica}
\label{cap:fundamentacao}

Este capítulo articula três corpos teóricos que sustentam o restante do trabalho: o orçamento público brasileiro como objeto, o controle social como propósito e o Data Warehouse como meio. Os três se cruzam no argumento central: a forma como os dados públicos são estruturados condiciona quem consegue interpretá-los.

\section{Orçamento Público e Transparência Governamental}

O orçamento público brasileiro se organiza em três instrumentos integrados: o Plano Plurianual (PPA), a Lei de Diretrizes Orçamentárias (LDO) e a Lei Orçamentária Anual (LOA) \cite{giacomoni2021orcamento}. O PPA define metas e diretrizes para um período de 4 anos. A LDO estabelece prioridades para o exercício seguinte. A LOA estima receitas e fixa despesas com base anual. Esse encadeamento confere ao orçamento brasileiro uma característica importante para a análise: as decisões financeiras de um ano específico não nascem isoladas, mas decorrem da execução de planos plurianuais.

As despesas são classificadas em quatro dimensões: institucional (qual órgão executa), funcional (a que função de governo se destina), programática (qual programa atende) e por natureza (tipo de gasto). Cada despesa percorre três estágios distintos antes de virar pagamento efetivo: o empenho reserva o crédito orçamentário, a liquidação verifica o direito do credor e o pagamento efetiva a transferência. A separação entre essas etapas explica fenômenos comuns no debate público, como os "restos a pagar" e as diferenças entre o valor empenhado e o valor pago. Sem entender essa cadeia, é fácil tirar conclusões equivocadas a partir dos dados.

Documentos normativos, como o Manual de Contabilidade Aplicada ao Setor Público (MCASP) e o Manual Técnico de Orçamento (MTO), padronizam essas classificações em âmbito nacional \cite{MCASP,MTO}. A padronização é uma virtude para o controle interno e para a contabilidade pública, mas se torna um obstáculo quando o destinatário é o cidadão. Pesquisas indicam que termos como "elemento de despesa", "função programática" e "fonte de recurso" atuam como barreiras semânticas e dificultam a interpretação por não especialistas \cite{melo2019proposta,marco2022transparencia}. O problema não é a falta de padrão, mas sim um padrão pensado para o emissor, e não para o receptor.

\section{Controle Social e Transparência}

A LAI consolidou a transparência governamental como obrigação legal, mas estudos consistentes mostram que publicar não é o mesmo que prestar contas. \citeonline{pinho2009accountability} discute essa distância entre transparência formal e \textit{accountability} efetiva no contexto brasileiro: a primeira tem caráter declarativo, a segunda exige resposta, justificativa e eventual sanção. Há um abismo entre disponibilizar um arquivo CSV no portal e responder à pergunta que o cidadão de fato faria. Em registro próximo, \citeonline{baldo2018democratizacao} sustenta que a democratização do orçamento exige a articulação simultânea de três princípios, a saber, legalidade, legitimidade e economicidade. Privilegiar qualquer um deles, isoladamente, compromete o resultado: legalismo burocrático sem participação real, economicismo que entrega o espaço público a interesses privados ou legitimidade sem ancoragem que gera expectativas inviáveis.

A participação cidadã qualificada esbarra na linguagem técnico-burocrática. \citeonline{marco2022transparencia} mostram, em estudo com Observatórios Sociais, que mesmo organizações civis dedicadas à fiscalização enfrentam dificuldade para cruzar dados de portais municipais. Se grupos organizados encontram esse obstáculo, o cidadão isolado dificilmente terá condições melhores.

Paulo Freire oferece um vocabulário próprio para discutir esse problema. Sua noção de conscientização aponta para a capacidade de "ler o mundo" criticamente, indo além da decodificação passiva de signos \cite{freire1967liberdade,freire1996autonomia}. \citeonline{streck2017educar} investigam empiricamente a presença desse legado em experiências de orçamento participativo no Rio Grande do Sul e mostram que as formulações inspiradas em Freire convivem com fatores que limitam seu potencial formativo. A intenção pedagógica das políticas participativas, na prática, costuma se descolar do que de fato chega ao cidadão. Aplicada ao orçamento público, a ideia ganha contornos técnicos. A conscientização orçamentária implica acessar os dados, compreender o significado dos termos, identificar padrões ao longo do tempo e formular perguntas próprias. A organização da informação não é neutra. Ela define o universo de leitores possíveis e, com isso, o universo de fiscalizadores possíveis.

\section{Data Warehouse e Modelagem Dimensional}

O Data Warehouse (DW) é um repositório voltado à análise, e \citeonline{kimball2013data} resumem suas propriedades em uma definição clássica: orientado a assunto, integrado, não volátil e variável no tempo. Cada propriedade tem efeito prático no uso analítico.

Orientado a assunto significa que os dados são organizados em torno de temas relevantes para análise, como despesas e receitas, e não em torno de processos transacionais como "cadastro de empenho". A diferença entre as duas organizações se torna evidente quando o analista pergunta "quanto a saúde gastou em 2024". Em um sistema transacional, essa pergunta exige a junção de dezenas de tabelas. Em um DW, ela cabe em uma consulta de poucas linhas.

Integrado quer dizer que dados de fontes heterogêneas chegam ao DW sob convenções uniformes de nome, tipo, unidade e formato. No caso público, isso resolve um problema concreto: arquivos de anos diferentes do mesmo portal costumam usar nomenclaturas distintas para o mesmo conceito. A integração elimina o esforço de reconciliação manual que, hoje, qualquer usuário sério dos portais é obrigado a realizar.

Não volátil indica que os dados carregados não sofrem atualizações destrutivas, apenas se acrescentam novos dados. Essa propriedade preserva o histórico e confere ao DW a característica de arquivo público auditável, em que nenhuma versão antiga é substituída silenciosamente.

Variável no tempo significa que o histórico é tratado como uma dimensão analítica. Cada fato carrega sua marcação temporal. Isso permite análises que comparam anos, identificam sazonalidades e detectam mudanças de padrão. Para a execução orçamentária, é justamente aí que residem as perguntas mais relevantes.

A modelagem dimensional, consolidada por Kimball e Ross, organiza o DW em torno do esquema em estrela. Tabelas de fatos armazenam métricas quantitativas (no caso orçamentário, valor empenhado, valor liquidado e valor pago). Tabelas de dimensão fornecem contexto descritivo: tempo, órgão, função, programa, natureza da despesa, favorecido, fonte de recurso. Essa arquitetura prioriza a usabilidade analítica e dispensa o usuário de dominar as junções complexas típicas de modelos relacionais normalizados.

Quando combinados com \textit{dashboards} interativos, os DW transformam milhares de linhas em visualizações navegáveis. O que estava disperso em planilhas vira um gráfico temporal, um mapa de calor e um comparativo entre áreas. Essa mediação tecnológica entre o dado bruto e o leitor leigo é o ponto de encontro entre as teorias de Kimball e de Freire. Kimball oferece a estrutura e Freire, o propósito. Esse trabalho parte da hipótese de que ambos precisam andar juntos no projeto de uma transparência efetiva.
